\section{Discussion}
\label{sec:discussion}

\paragraph{What the iterated local update buys.}
The architecture comparison in \S\ref{sec:results-baselines}
isolates one question: holding the conditioner constant, does the
NCA's iterated-local-update body outperform a non-iterative body?
Our findings suggest the answer is yes on tasks where the engine's
\texttt{again} loop produces non-trivial cascade chains, and roughly
a wash on tasks where one tick of dynamics resolves in a single rule
firing. The NCA's parameter efficiency is also striking: shared
weights with $T \cdot n_\ell$ effective depth use the same parameters
as a one-step model.

\paragraph{Memorization as a confound.}
A key observation in \S\ref{sec:results-ceiling} is that the
single-level / single-grid setting allows the model to learn a
per-state lookup table via cross-attention over rich slot
representations. ``Pool features substitute for depth'' is the
single-level reading; ``the model memorizes per-state outcomes'' is
the multi-level reading of the same architectural configuration.
Both are real, and disentangling them required the multi-grid
synthetic-data regime that we describe in \S\ref{sec:results-ceiling}.
We recommend the multi-grid synth setup as a default check for
future PuzzleScript world-model work: if your model only succeeds at
the single-grid setting, you may not be learning the dynamics.

\paragraph{Open: escaping the multi-grid ceiling.}
Our experiments characterize the ceiling but do not break it. The
remaining candidate fixes that have not been ruled out are
(a) hard / sparse rule-slot routing so the model cannot smear
``fire-once'' across all slots, (b) per-cell halt mechanism so cells
in their fixed point stop updating, (c) curriculum or replay-buffer
training that selectively exposes the model to slide-distance
$\geq 2$ instances, and (d) auxiliary ``is-changing'' mask head
factoring ``where to fire'' from ``what to write.'' A successful
break here would close the most important remaining gap between
authored-data and synthetic-multi-grid performance.

\paragraph{Beyond PuzzleScript.}
The architecture is not domain-specific. Any deterministic
discrete-grid simulator with declarative rules (Conway's Game of
Life and its variants, cellular-automata gas dynamics, Asari-style
chemical-reaction grids) admits the same modeling pattern: tokenize
the rules, encode them via the slot encoder, train an NCA body
shared across iterations. We expect the comparison against
non-iterative baselines to favor the NCA more strongly as the
underlying engine's outer-loop iteration count grows.

\paragraph{Limitations.}
Our slot encoder is level-independent: the slots depend only on the
game's rule tokens, not on the level's spatial state. This is the
intended design (one encoded game $\Rightarrow$ many levels) but it
also caps what the slots can carry. A version where slots can attend
back into the input grid (perceiver-IO style) would be a natural
extension. We also restrict to single-step prediction; multi-step
training (predict $s_{t+k}$ for $k > 1$) is straightforward to add
and not yet explored.
